%==============================================================================
%  CAPÍTULO — EL MODELO ESPAÑOL Y EL MARCO EUROPEO
%==============================================================================
\chapter{El modelo español y el marco europeo}
\label{cap:comparado-espana}

\fichaCapitulo{El comparador más cercano al derecho chileno por tradición jurídica.}{Texto
Refundido de la Ley Concursal, reforma de la Ley 16/2022, exoneración del pasivo insatisfecho y
la Directiva (UE) 2019/1023 sobre reestructuración preventiva.}

%==============================================================================
\bloqueTeorico
%==============================================================================

\subsection{Por qué España es el comparador natural}

Si el derecho estadounidense aporta los modelos institucionales, el español aporta algo distinto y
para el jurista chileno igualmente valioso: \destacar{un ordenamiento de la misma familia
jurídica}, con idéntica tradición de derecho civil codificado, garantías reales reguladas en
términos equivalentes y una dogmática construida sobre las mismas categorías ---prelación de
créditos, acción pauliana, par condicio---. Los trasplantes desde España no exigen la traducción
conceptual que sí requiere el derecho angloamericano.

A ello se suma una circunstancia práctica: buena parte de los casos chilenos de insolvencia
transfronteriza han tenido su procedimiento principal en España ---Elimco, Isolux, Assignia---, de
modo que conocer el concurso español es, para el concursalista chileno, un requerimiento operativo
y no sólo académico (capítulo \ref{cap:transfronteriza}).

\subsection{Evolución: de la Ley Concursal de 2003 al Texto Refundido}

El derecho español vivió, en dos décadas, un ciclo de reformas que Chile recorrió de manera
comprimida. La Ley Concursal de 2003 unificó los antiguos procedimientos ---quiebra, suspensión de
pagos, concurso de acreedores y quita y espera--- en un procedimiento único, del mismo modo que la
\LIR{} sustituyó en Chile el régimen de la Ley \Nro 18.175. Sucesivas reformas incorporaron los
institutos preconcursales, y el \destacar{Texto Refundido de la Ley Concursal} de 2020 (TRLC)
sistematizó el conjunto.

El hito decisivo es la \destacar{Ley 16/2022} \parencite{esp-ley16-2022}, que transpuso la
Directiva europea sobre reestructuración e insolvencia y reconfiguró tanto el instrumento
preventivo como la segunda oportunidad de la persona natural.

\subsection{La Directiva (UE) 2019/1023 como marco supranacional}

La \destacar{Directiva (UE) 2019/1023} \parencite{ue-directiva-2019-1023} sobre marcos de
reestructuración preventiva, exoneración de deudas e inhabilitaciones constituye hoy el estándar
mínimo europeo. Sus tres pilares:

\begin{enumerate}
  \item \textbf{Marcos de reestructuración preventiva.} Los Estados deben ofrecer al deudor con
  \emph{probabilidad de insolvencia} ---no insolvencia actual--- un procedimiento para
  reestructurarse antes del colapso, con suspensión de ejecuciones singulares por un plazo inicial
  limitado (típicamente cuatro meses, prorrogable hasta doce).

  \item \textbf{Exoneración plena de deudas.} La persona natural insolvente debe poder acceder a
  una descarga total en un plazo \destacar{no superior a tres años}, con posibilidad de que los
  Estados excluyan ciertas categorías de crédito.

  \item \textbf{Reestructuración con arrastre de clases.} Se admite la confirmación judicial del
  plan pese a la oposición de una o más clases ---el \emph{cross-class cram-down}---, sujeta a
  contrapesos de equidad.
\end{enumerate}

\begin{foco}[La probabilidad de insolvencia como umbral de acceso]
El aporte conceptual más importante de la Directiva es adelantar el umbral: no exige cesación de
pagos ni insolvencia actual, sino \destacar{probabilidad de insolvencia}. La lógica es que la
reestructuración temprana conserva valor y la tardía sólo administra la pérdida. La \LIR{} chilena
no fija un umbral explícito para la reorganización, lo que en la práctica permite un acceso
temprano, pero tampoco lo incentiva mediante señales normativas claras.
\end{foco}

%==============================================================================
\bloqueNormativo
%==============================================================================

\subsection{Los planes de reestructuración y el arrastre de clases}

La reforma española de 2022 sustituyó los antiguos acuerdos de refinanciación por los
\destacar{planes de reestructuración}, cuyo rasgo distintivo es la posibilidad de imponerse a
clases disidentes. El esquema es el siguiente:

\paragraph{Formación de clases.} Los créditos se agrupan en clases según un
\destacar{interés común}, atendiendo a su rango concursal. La formación de clases es impugnable, y
constituye ---como en el Capítulo 11 estadounidense--- el terreno donde se libra buena parte de la
batalla del plan: quién queda con quién determina las mayorías.

\paragraph{Arrastre dentro de la clase.} Aprobado el plan por la mayoría de una clase, vincula a
los disidentes de esa misma clase.

\paragraph{Arrastre entre clases (\emph{cross-class cram-down}).} El plan puede homologarse pese al
rechazo de clases enteras, si lo aprueba una mayoría de clases ---o al menos una clase que estaría
\emph{in the money} en una liquidación--- y se cumplen dos garantías:

\begin{itemize}
  \item La \destacar{prueba del interés superior de los acreedores}: ningún disidente puede recibir
  menos de lo que obtendría en una liquidación concursal.
  \item La \destacar{regla de prioridad}, en su versión absoluta o en la variante \emph{relativa}
  que la Directiva permite, y que España adoptó de forma matizada: las clases superiores deben
  recibir un trato al menos tan favorable como las inferiores.
\end{itemize}

\begin{alerta}[El contraste con el artículo 96 chileno]
En Chile, el rechazo de la propuesta por una clase conduce derechamente a la liquidación
(\art{96}), sin que exista mecanismo alguno para superar la oposición de una minoría de clase.
El sistema chileno confía enteramente en la negociación previa; el europeo admite que el juez
imponga el plan cuando la oposición resulta económicamente irracional a la luz de la alternativa
liquidatoria. Es la asimetría más relevante entre ambos modelos y la que explica la discusión sobre
las quitas al acreedor garantizado disidente examinada en el capítulo
\ref{cap:reorganizacion-ordinaria}.
\end{alerta}

\subsection{La exoneración del pasivo insatisfecho (EPI)}

La segunda oportunidad española es el comparador directo del \discharge{} chileno del \art{255}. La
Ley 16/2022 la reconfiguró en términos que conviene precisar:

\paragraph{Dos modalidades.} El deudor puede optar por la exoneración
\destacar{con plan de pagos} ---conservando bienes, incluida eventualmente la vivienda, a cambio de
cumplir un plan de hasta tres o cinco años--- o por la exoneración
\destacar{con liquidación previa} de su patrimonio. La primera modalidad, que permite conservar la
vivienda habitual, carece por completo de equivalente chileno.

\paragraph{El requisito de buena fe.} La exoneración exige un deudor de buena fe, noción que la ley
define \destacar{negativamente}, mediante una lista tasada de conductas que la excluyen ---condena
por delitos económicos, incumplimiento de deberes de colaboración, información falsa, entre
otras---. Esta técnica legislativa es exactamente la del \art{169 A} chileno, que también tasa las
causales de mala fe en lugar de definir la buena fe en positivo.

\paragraph{El crédito público.} La ley excluye de la exoneración el crédito público más allá de
umbrales determinados ---del orden de \destacar{diez mil euros} por cada uno de los créditos
tributario y de seguridad social---.\verificar{Confirmar los umbrales vigentes del crédito público
exonerable en la Ley 16/2022 tras sus modificaciones posteriores.}

\begin{practica}[El paralelo exacto con la controversia chilena del CAE]
La exclusión española del crédito público es la misma cuestión de fondo que el CAE chileno: hasta
dónde puede el interés fiscal resistir el efecto liberatorio del concurso. La diferencia es de
técnica y de transparencia. España lo resolvió \destacar{en la ley}, con umbrales cuantitativos
discutibles pero explícitos y sujetos a debate parlamentario. Chile lo dejó a la
\destacar{construcción judicial} por vía del principio de especialidad, con el resultado de
sentencias contradictorias entre salas. Este contraste es, probablemente, el argumento comparado
más fuerte a favor de una reforma legal que fije un catálogo chileno de deudas no exonerables.
\end{practica}

\subsection{La administración concursal y los auxiliares}

El \destacar{administrador concursal} español concentra funciones que en Chile se reparten entre el
veedor y el liquidador. Su designación corresponde al juez del concurso, con criterios de turno y
especialización, y su estatuto de responsabilidad, retribución arancelada y rendición de cuentas
guarda notable paralelismo con el régimen chileno de los \art{35} y siguientes.

Dos diferencias merecen registro. La retribución española se fija mediante \destacar{arancel
reglamentario} sobre el activo y el pasivo, y no mediante una tabla legal progresiva sobre lo
efectivamente repartido como el \art{40} chileno ---lo que altera los incentivos: el sistema
chileno premia el recupero efectivo, el español remunera la magnitud del concurso---. Y el juez
español conserva un rol de dirección considerablemente más activo que el del juez civil chileno.

%==============================================================================
\bloquePractico
%==============================================================================

\subsection{Tabla de equivalencias funcionales}

\begin{table}[htbp]
\centering
\small
\caption{Correspondencias entre el derecho español y la Ley \Nro 20.720}
\label{tab:cap26-equivalencias}
\begin{tabularx}{\textwidth}{@{}>{\raggedright\arraybackslash}p{0.28\textwidth}
                                 >{\raggedright\arraybackslash}p{0.26\textwidth} X@{}}
\toprule
\textbf{España} & \textbf{Chile} & \textbf{Diferencia relevante} \\
\midrule
Planes de reestructuración & Acuerdo de reorganización judicial & España admite arrastre entre
clases; Chile no \\
\addlinespace
Formación de clases impugnable & Clases del \art{61} & En España la formación de clases es en sí
misma objeto de impugnación \\
\addlinespace
Prueba del interés superior de los acreedores & \emph{Sin equivalente expreso} & Piso de
liquidación como garantía del disidente \\
\addlinespace
Exoneración con plan de pagos & \emph{Sin equivalente} & Permite conservar la vivienda habitual \\
\addlinespace
Exoneración con liquidación previa & Liquidación + \art{255} & Modelo funcionalmente equivalente \\
\addlinespace
Buena fe definida por lista negativa & Causales tasadas del \art{169 A} & Misma técnica
legislativa \\
\addlinespace
Exclusión legal del crédito público & Exclusión judicial del CAE & España lo resuelve en la ley;
Chile, por vía jurisprudencial \\
\addlinespace
Administrador concursal único & Veedor y liquidador & Chile separa las funciones según el
procedimiento \\
\addlinespace
Retribución por arancel & Tabla del \art{40} sobre repartos & Chile incentiva el recupero
efectivo \\
\bottomrule
\end{tabularx}
\end{table}

\begin{foco}[Tres lecciones del modelo europeo para Chile]
Primera: \destacar{legislar las exclusiones del \discharge{}} en lugar de dejarlas a la
construcción judicial ---la solución española, aun con umbrales opinables, es preferible a la
incertidumbre chilena---. Segunda: incorporar el \destacar{arrastre entre clases} con los
contrapesos de la Directiva, que convertiría el acuerdo de reorganización en un instrumento
utilizable frente a minorías obstruccionistas. Tercera: contemplar una modalidad de exoneración
\destacar{con plan de pagos} que permita a la persona deudora conservar la vivienda, hoy ausente en
el catálogo chileno y probablemente la carencia más sensible del sistema para el deudor
sobreendeudado con inmueble.
\end{foco}
